\documentclass[12pt]{article}
\usepackage[margin=0.6in,top=1.15in,bottom=0.55in]{geometry}
\usepackage{lmodern}
\usepackage{microtype}
\usepackage{titlesec}
\usepackage{enumitem}
\usepackage[hidelinks]{hyperref}
\usepackage{../../latex/grader-worksheet}

\setlength{\parindent}{0pt}
\setlength{\parskip}{0.25em}
\titleformat{\section}{\large\bfseries\scshape}{}{0pt}{}

\GraderSetup{assignment-id=U1L15LL,template-revision=1}

\begin{document}

\begin{center}
{\LARGE\bfseries Week 15 Logic Lab}\par\vspace{0.05em}
{\large\scshape Verbal and Real Disputes}\par\vspace{0.1em}
\rule{0.78\textwidth}{0.5pt}\par\vspace{0.15em}
\end{center}

\noindent\textbf{Purpose:} Apply Week 15 skills to new disagreements: clarify terms, identify shared claims, and classify disputes as verbal, real, or mixed.

\vspace{0.3em}

\section*{Part 1: Find the Contested Term}

\textbf{Question 1.1}\\[-0.15em]
Two students argue, ``The project is finished.'' One means every slide exists; the other means the sources are checked and the speaking notes are practiced. Identify the contested term and each side's meaning.

\GraderAnswerBox{Part 1 Q1}{2.35in}

\textbf{Question 1.2}\\[-0.15em]
A family argues whether a weekend trip is ``expensive.'' One person means the total cost is high; another means the cost is too high for this month's budget. Identify the contested term and each side's meaning.

\GraderAnswerBox{Part 1 Q2}{2.35in}

\newpage

\section*{Part 2: Verbal, Real, or Mixed?}

\textbf{Question 2.1}\\[-0.15em]
A coach says, ``Practice was successful because everyone attended.'' A captain says, ``Practice was not successful because the team did not improve the weak plays.'' Classify the dispute and explain what clarification is needed.

\GraderAnswerBox{Part 2 Q1}{2.55in}

\textbf{Question 2.2}\\[-0.15em]
Two voters agree that a safe street means fewer crashes and fewer serious injuries. One says the new speed bumps made the street safer; the other says they did not. Classify the dispute and name what evidence would now matter.

\GraderAnswerBox{Part 2 Q2}{2.45in}

\newpage

\section*{Part 3: Shared Claim and Remaining Disagreement}

\textbf{Question 3.1}\\[-0.15em]
A club debates whether a new member is ``committed.'' One side means the member attends every meeting; the other means the member completes assigned work between meetings. Name a shared claim that may emerge after clarification and the remaining disagreement.

\GraderAnswerBox{Part 3 Q1}{2.55in}

\textbf{Question 3.2}\\[-0.15em]
A class argues whether an essay is ``original.'' One side means it uses no copied sentences; the other means it makes a fresh argument rather than rearranging common points. Name what clarification may settle and what may still need judgment.

\GraderAnswerBox{Part 3 Q2}{2.55in}

\newpage

\section*{Part 4: Dispute Map in a Paragraph}

\textbf{Question 4.1}\\[-0.15em]
Read the paragraph. Make a dispute map: disputed claim, contested term, Side A meaning, Side B meaning, shared claim, remaining disagreement, and classification.

\vspace{0.2em}

\noindent\fbox{%
\begin{minipage}{\dimexpr\textwidth-2\fboxsep-2\fboxrule}
\vspace{0.35em}
The student council argues whether the lunch policy is ``equal.'' One group says it is equal because every student receives the same thirty-minute lunch period. Another group says it is not equal because students in the far science building lose ten minutes walking, while students near the cafeteria do not. Both groups agree that the written schedule gives everyone thirty minutes, and both agree that actual eating time matters.
\vspace{0.35em}
\end{minipage}
}

\GraderAnswerBox{Part 4 Q1}{3.2in}

\newpage

\section*{Part 5: Review}

\textbf{Question 5.1}\\[-0.15em]
``We proved the courtyard floods because several students imagined that the drain is clogged.'' Use a prior-week tool: what is the hypothesis, what must be discovered, and what would count as proof?

\GraderAnswerBox{Part 5 Q1}{2.2in}

\textbf{Question 5.2}\\[-0.15em]
``The new study app is reliable because three advertisements quote the same celebrity endorsement.'' Use a prior-week source test: what is weak about the corroboration, and what evidence would be more relevant?

\GraderAnswerBox{Part 5 Q2}{2.05in}

\end{document}
